\section*{Sources \& limitations}

本号では、公式発表・documentation・model card・プロジェクトのrelease、報道、コミュニティ上の反応を、それぞれの性質に合わせて読み分ける。公開情報が示す範囲と、まだ確かめられていない範囲を分けるためである。

\begin{tcolorbox}[enhanced,breakable,colback=SurveySoft,colframe=SurveyRule,boxrule=0.4pt,arc=1mm,left=3mm,right=3mm,top=2mm,bottom=2mm]
\textbf{公式・プロジェクト資料}\par
モデルやソフトウェアの提供条件、利用経路、機能、releaseの変更点を示す。提供元やプロジェクト自身が報告する速度・性能値は、独立測定の結果とは区別して読む。\par\smallskip
\textbf{報道}\par
報道が伝える仕様、時系列、測定値は、重み置き場や一次発表での確認を待つ段階では事実に格上げしない。出どころと留保ごと受け止める。\par\smallskip
\textbf{コミュニティ上の反応}\par
W35の観測では、公開の重み群と道具面への関心が見られた。これは話題の広がりを示すcontext-onlyの背景であり、技術性能や互換性の証明には用いない。\par\smallskip
\textbf{まだ確かめる必要がある点}\par
重みの利用条件、previewと旗艦の境界、利用対象とsafeguard、benchmarkの条件、異なるモデル・環境への一般化などは、公開情報が増えるまで断定しない。
\end{tcolorbox}

本文では、各資料の主体、日付、提供経路、測定条件をできるだけ明示した。異なる主体が報告した数値や、条件の違うbenchmarkを同じ尺度のランキングへ変換していない。

\nocite{glm53flash,qwen38flash,hy4preview,w35community,granite42,vscodeahp,cursororigin,copilotteams,harness,grith,fable5,gpt56kiro,vllm028,openthai20,autoalign,chinamiit,koreaethics}

\begin{claimboundary}[Temporal and claim boundary]
本文の基準時点は2026-08-28 18:00 America/New\_York。公式資料、プロジェクトのrelease、報道の伝える内容、コミュニティ上の観測は、それぞれ異なる範囲の主張を支える。提供元の速度値、報道の仕様、実験室の評価値、話題で見られた関心は相互に代替できず、測定条件が異なる値も単純比較できない。未解決の点は未解決のまま記載する。
\end{claimboundary}
